A model that maps the customer journey from first becoming aware of the product being sold all the way to them making a purchase (and into after-sales).

The name comes from the shape. For example, you start with a large pool of potential customers at the top, and a smaller number make it through each stage to eventually convert.

[insert diagram]

Typically, the stages move from:
\begin{enumerate}
    \item \textbf{Awareness.}
    \begin{enumerate}
        \item Prospects discover the brand via marketing (ads, social media, content). This comes from pursuing inbound marketing, where businesses attract customers by providing information that is useful or positioning themselves as visible.
    \end{enumerate}

    \item \textbf{Interest.}
    \begin{enumerate}
        \item Tools such as webinars, case studies, and landing pages aid in educating and engaging curious prospects.
        \item Often use educational tools to engage their prospects, including:
        \begin{enumerate}
            \item Webinars.
            \item Case studies.
            \item Landing pages.
            \item Informational content.
        \end{enumerate}
    \end{enumerate}

    \item \textbf{Consideration.}
    \begin{enumerate}
        \item Prospects evaluate the product meets their needs:
        \begin{enumerate}
            \item Pricing.
            \item Product features.
            \item NFQs.
            \item Customer Support.
            \item Trade-Offs with Competing Alternatives.
        \end{enumerate}
    \end{enumerate}

    \item \textbf{Intent.}
    \begin{enumerate}
        \item Clear purchase signals are any actions that a prospect takes that suggest they are moving toward a buying decision:
        \begin{enumerate}
            \item Requesting a product demo.
            \item Asking for a quote or pricing breakdown.
            \item Starting a free trial.
            \item Visiting pricing or checkout pages.
        \end{enumerate}
        \item Retargeting begins here. Retargeting is a marketing tactic where you serve ads to people who have not been converted yet. Usually work through tracking tools such as cookies.
    \end{enumerate}

    \item \textbf{Evaluation/Purchase.}
    \begin{enumerate}
        \item A thorough evaluation must be done, considering ROI, scalability, and integration with existing systems. During this stage, the \textbf{sales team plays a key role} by answering questions, providing product demonstrations, and addressing concerns. If the evaluation is successful, the prospect converts into a \textbf{paying customer}.
        \item Usually at this stage, prospects tend to finalize their purchase and become paying customers.
    \end{enumerate}

    \item \textbf{Loyalty (during After-Sales).}
    \begin{enumerate}
        \item Businesses focus on satisfaction and advocacy through excellent customer service. Ongoing engagement is important.
        \item Businesses focus on:
        \begin{enumerate}
            \item High-quality customer service.
            \item Continuous engagement.
            \item Customer satisfaction.
            \item Satisfied customers can become \textbf{advocates} who recommend the product to others.
        \end{enumerate}
    \end{enumerate}

\end{enumerate}

The following are some advantages to having a sales funnel:
\begin{enumerate}
    \item \textbf{Allows for better customer retention.}
    \begin{enumerate}
        \item Acquiring a new customer is significantly more expensive than keeping one. Retained customers also tend to spend more overtime.
    \end{enumerate}

    \item \textbf{Using real behavioral data, clicks, purchases, drop-off points, to make smarter decisions instead of guessing.}
    \begin{enumerate}
        \item Without data, optimizing is much harder. With data, you can see exactly where prospects are getting stuck or dropping off.
        \item CRM such as SalesForce, or HubSpot, log every interaction a prospect has with the brand. The funnel feeds data into the system, giving your team a full picture of each customer's history and behavior.
    \end{enumerate}

    \item \textbf{Enhanced audience targeting.}
    \begin{enumerate}
        \item Marketing resources, like all resources, are scarce, and finite. Spending them on bad leads to waste energy.
        \item The funnel helps build buyer personas, profiles of the ideal customer.
    \end{enumerate}

    \item \textbf{Increased referrals and social proof.}
    \begin{enumerate}
        \item People trust people more than they trust brands (not always true for established brands, but for new brands this tends to be the case).
        \item By delivering good experience, especially after-sales, advocates for the product can be formed.
    \end{enumerate}

\end{enumerate}

\subsubsection{Growing Role of AI in Sales Funnel}

IBM did a study through IBMV that found that by 2026, 83\% of executives expect AI agents to automatically execute real-time tasks and automated deal progression. Embracing AI must be strategic, and now there is more of an emphasis on cleaning data within the ETL pipeline. (not sure where this fits in).

AI systems can automatically score leads, personalize marketing messages, trigger follow-ups at optimal times, and predict which prospects are most likely to convert.